\section{两类项目路径：集成 Existing Benchmark 或构建 New Benchmark}

slide 72 将课程项目分为两类。Type 1 选择已发表 benchmark，完成 AgentBeats 集成并分析原 benchmark 的质量问题；Type 2 从新能力出发构造 benchmark。两条路径都不是简单软件作业：前者必须研究 validity、bias 和 limitation，后者必须证明任务价值、环境真实性与评分可靠。

\lecturefigure{slides-images/slide-072.png}{官方 slide 72：Existing benchmark integration 与 new benchmark 两类项目}

\subsection{Type 1：集成不等于包装 API}

slides 73--74 规定 Type 1 的目标和额外分析。团队要阅读论文/代码，明确 task formulation、数据来源、metric 与已知限制；把原 interface 适配为平台协议；复现实验；再做 benchmark quality analysis，例如 contamination、难度、bias、reliability、environment difference 和 failure mode。

\lecturefigure{slides-images/slide-073.png}{官方 slide 73：集成现有 benchmark 的目标、团队与参考材料}
\lecturefigure{slides-images/slide-074.png}{官方 slide 74：除平台集成外还需 benchmark quality analysis}

\begin{importantbox}{复现先于改造}
先在原仓库/原协议下复现论文基线，再做 Green Agent 适配。否则分数差异无法判断来自 participant、环境版本、接口实现还是 benchmark 本身。保存原始 commit、依赖锁、数据 checksum 和 baseline output，形成迁移前后的证据链。
\end{importantbox}

\subsection{Type 1 workflow 与 SWE-bench Verified}

slide 77 是 progressive build 的最终工作流：integration、benchmark quality analysis、metrics、bias/limitation、correction/expansion。slide 78 用 SWE-bench 与 Verified 说明已有 benchmark 也可能包含难以复现、测试不公平或题目不清的问题；人工筛选和修订可提高 outcome validity，但也改变任务分布，需要公开版本差异。

\lecturefigure{slides-images/slide-077.png}{官方 slide 77：Type 1 的集成、质量分析、指标与修订完整 workflow}
\lecturefigure{slides-images/slide-078.png}{官方 slide 78：SWE-bench 与 SWE-bench Verified 的问题筛查}

\begin{knowledgebox}{读图：Verified 的价值和边界}
Verified 集合通过人工审查提升 task clarity 和 test fairness，使分数更可解释。它不意味着原集合所有未入选任务都无价值，也不证明新集合没有污染或覆盖偏差。报告应注明 benchmark 版本，并分别讨论 reliability 与 representativeness：更可靠的小集合可能覆盖更窄。
\end{knowledgebox}

\begin{warningbox}{Correction 不能偷偷改标尺}
修复测试、环境或题目后，历史分数通常不可直接比较。应发布新版本、迁移说明、受影响 task 列表和 baseline rerun。若只在本地修补却沿用原 benchmark 名称，会制造无法审计的 leaderboard。
\end{warningbox}

\subsection{Type 2：新 benchmark 从有用场景和可验证 outcome 出发}

slides 79--80 要求新 benchmark 覆盖尚未被良好测量的能力，任务应反映 useful real-world scenario，并能持续、低成本生成或维护。数据来源可以是人工、合成、真实日志或其组合；关键是难度、真实性、许可和 oracle。团队还需说明为何现有 benchmark 不足，而不是仅换领域名词。

\lecturefigure{slides-images/slide-079.png}{官方 slide 79：从零创建 benchmark 的目标}
\lecturefigure{slides-images/slide-080.png}{官方 slide 80：真实场景、持续数据生成与任务来源要求}

\begin{knowledgebox}{新 benchmark 的 novelty 检查}
能力新颖：测到现有套件缺失的行为；环境新颖：引入真实工具/状态约束；测量新颖：更可靠地验证已有任务；数据新颖：覆盖新的分布或动态生成。Novelty 应对应评估价值，而不是为了不同而不同。若只是换数据集，必须证明分布变化会暴露新的失败模式。
\end{knowledgebox}

\begin{importantbox}{可扩展任务生成的质量门}
任务生成器应输出 task、initial state、allowed tools、ground truth/constraints 和 provenance；随后经过 schema validation、可执行性检查、难度估计、去重、污染扫描和人工抽样。生成速度不是目标，能持续产生有效、可分离且不泄漏答案的任务才是。
\end{importantbox}

\subsection{本章小结}

Type 1 强调可靠迁移与质量审计，Type 2 强调任务价值与测量设计。两者都要版本化环境和指标，并用 baseline 证明 Green Agent 的实现没有改变原本要测的能力。

